\begin{table}[htbp]
\centering
\caption{Summary Statistics: State-Level Consumer Price Index (General)}
\label{tab:summary}
\begin{tabular}{lcccccc}
\toprule
& \multicolumn{2}{c}{Full Sample} & \multicolumn{2}{c}{Pre-GST} & \multicolumn{2}{c}{Post-GST} \\
\cmidrule(lr){2-3} \cmidrule(lr){4-5} \cmidrule(lr){6-7}
Variable & Mean & SD & Mean & SD & Mean & SD \\
\midrule
CPI Index & 150.5 & 28.2 & 121.2 & 9.1 & 166.3 & 21.6 \\
Log CPI Index & 4.996 & 0.187 & 4.795 & 0.076 & 5.105 & 0.13 \\
Inflation (\%) & 4.9 & 2.6 & 5.5 & 2.8 & 4.7 & 2.4 \\
Tax Intensity (raw \%) & 5.6 & 1.6 & 5.6 & 1.6 & 5.6 & 1.6 \\
Tax Intensity (std) & 0 & 0.986 & 0 & 0.986 & 0.001 & 0.986 \\
\midrule
\multicolumn{7}{l}{States: 35 \quad Months: 154 (54 pre-GST, 100 post-GST) \quad Observations: 5,387} \\
\bottomrule
\end{tabular}
\begin{tablenotes}[flushleft]\small
\item \textit{Notes:} CPI base year 2012=100 (Combined sector). Tax intensity is pre-GST state indirect tax revenue as \% of GSDP (RBI State Finances 2016-17). Standardized intensity has mean 0 and SD 1. Pre-GST: January 2013--June 2017 (54 months). Post-GST: July 2017--December 2025 (100 months, excluding April--May 2020). Two months missing due to COVID-19 lockdown suspension of CPI collection.
\end{tablenotes}
\end{table}
